\section{The Absorption Problem}

\subsection{Definition}

We define the \textbf{Absorption problem} as the case where a model maps
multiple distinct epistemic states to the same scalar T/I/F vector,
typically (T=0, I$\approx$1, F=0). Under Absorption, indeterminacy absorbs the
truth and falsity dimensions, leaving the scalar output unable to
distinguish between types of uncertainty.

\subsection{Evidence}

Mistral Medium 3.1 predominantly exhibits Absorption across multiple
phenomena in S1 (modal response across 5 reps shown):

\begin{longtable}[]{@{}llllll@{}}
\toprule\noalign{}
Phenomenon & T & I & F & Sum & Consistency \\
\midrule\noalign{}
\endhead
\bottomrule\noalign{}
\endlastfoot
Paradox (Logical) & 0.00 & 1.00 & 0.00 & 1.00 & 4/5 reps \\
Ignorance (Epistemic) & 0.00 & 1.00 & 0.00 & 1.00 & 3/5 reps \\
Contingency (Future) & 0.32 & 0.71 & 0.11 & 1.14 & mean values \\
\end{longtable}

The original study's flagship model (GPT-4o) shows the same pattern:
(T=0, I=1, F=0) for both paradox and ignorance. Neither Leyva-Vázquez
nor Smarandache discuss this.

Absorption is not a model deficiency: it is a representational
limitation. The model may have richer internal distinctions that the
scalar output format cannot express. Testing this hypothesis motivates
the tensor extension.
